\section{Conclusion}\label{sec:conclusion}

In the housing subsidy scenario, the metamodel revealed that no credential format simultaneously satisfies the eIDAS ARF format mandate, GDPR data minimization, and W3C VCDM 2.0 conformance on the income credential: three governance frameworks, each internally consistent, whose joint requirements are unsatisfiable. The floor area constraint further exposed that cross-credential predicate evaluation across two independently issued credentials exceeds the capabilities of every deployed format.

The three-layer metamodel (DCL, \ac{CSL}, \ac{FSL}), grounded in \ac{VCDM} 2.0 and formalized as graph predicates in Refinery (\autoref{sec:approach}), makes cross-layer constraints from heterogeneous governance frameworks jointly evaluable. Validation confirmed coverage against the W3C specification, expressiveness against regulatory and standards sources, and error visibility against known design anti-patterns (\autoref{sec:evaluation}). Both headline results and two of five anti-pattern categories require predicates that reference elements from independently governed layers; no single-layer formalization can express them without collapsing the governance-source distinction that makes the constraints meaningful.

The primary limitation is the \ac{FSL}'s current scope: it models format capabilities as predicates over the class hierarchy but does not yet carry format-internal structural constraints (e.g., SD-JWT-VC disclosure granularity, mdoc namespace partitioning). Future work targets FSL format extensions, broader governance catalogs (national eIDAS implementations, sector-specific rules), practitioner evaluation, evaluation against alternative solvers (Alloy, USE/OCL), and range proof semantics in capability predicates.
